% MT030 / explanation — canonical editable content.
L1、L2正则化都在训练损失外增加对\textbf{权重}的惩罚。两者的共同作用是限制复杂度，具体优化方式与是否产生零系数却不同。



\begin{itemize}

\item \textbf{A不选。} L2在某些问题中改善求解稳定性，但L1在0处不可导，不能一概称两者都让优化更容易。

\item \textbf{B不选。} “良态”是指输入数据略有变化时，求出的系数不会剧烈变化。若一些特征几乎重复，普通最小二乘可能很不稳定。把每条输入排成一行得到数据矩阵\(X\)；Ridge求解时将\(X^TX\)换成\(X^TX+\alpha I\)，其中\(I\)的主对角线为1、其他位置为0。\(\alpha>0\)相当于给主对角线加一个正数，可以改善求解稳定性。L1没有同样的矩阵求逆形式，不能把Ridge的理由直接套过去。

\item \textbf{C不选。} 若数据损失仍是平方误差，一个极端输出值仍能带来很大的损失与梯度。限制权重与采用稳健残差损失是不同操作；绝对损失的作用见\StudyRef{Answers}{mt025}{MT025}。

\item \textbf{D按原解选，但拆开理解。} 两者都可以缩小系数；L1还能把某些系数直接变为0，自动选掉特征。Ridge通常只是让系数较小，若据此删特征还需要人为阈值等额外规则。

\item \textbf{E选。} 给权重增加代价，可减少为贴合训练细节而使用复杂参数的倾向，有助于控制过拟合；过强惩罚又可能欠拟合。

\end{itemize}


\textbf{出处：} 2021B*期中 Q4，原卷第2页（5分）


\textbf{答案依据：} 2021B*期中 Q4 随卷／配套参考答案；条件与解释经复核。
